% Poema grothendieckiano.
% Nota breve, alla maniera degli Éléments (sito, campo, gerbo, discesa, punto di vista relativo).
% Il vocabolario resta astratto di proposito: la natura dei due problemi si indovina leggendo.
\documentclass[10pt]{article}
\usepackage[a4paper,left=42mm,right=42mm,top=24mm,bottom=26mm]{geometry}
\usepackage{fontspec}
\usepackage{amsmath,amssymb,amsthm}
\usepackage[italian]{babel}
\usepackage[hidelinks]{hyperref}
\usepackage{enumitem}
\setlist{nosep,leftmargin=1.4em}
\setlength{\parskip}{0.25em}
\setlength{\emergencystretch}{2em}

\newtheoremstyle{bbk}{0.45em}{0.45em}{\itshape}{}{\bfseries}{.}{0.4em}{}
\theoremstyle{bbk}
\newtheorem{prop}{Proposizione}[section]
\newtheorem{coro}[prop]{Corollario}
\newtheorem{conj}[prop]{Congettura}
\newtheoremstyle{bbkdef}{0.45em}{0.45em}{}{}{\bfseries}{.}{0.4em}{}
\theoremstyle{bbkdef}
\newtheorem{defi}[prop]{Definizione}
\newtheorem{rema}[prop]{Osservazione}
\renewcommand{\qedsymbol}{\(\square\)}
\renewcommand{\proofname}{Dimostrazione}

\newcommand{\Lec}{\mathcal L}
\newcommand{\U}{\mathfrak U}
\newcommand{\Cl}[1]{\overline{#1}}
% Un mot caviardé : un rectangle noir de sa largeur ; le mot lui-même n'entre pas dans le PDF.
\newcommand{\cache}[1]{\setbox0\hbox{#1}\rule[-0.32ex]{\wd0}{1.02em}}
\makeatletter
\renewcommand{\@seccntformat}[1]{\S\,\csname the#1\endcsname.\quad}
\makeatother

\title{\vspace{-1.2em}\large\bfseries Poema grothendieckiano\\[0.2em]
\normalsize\mdseries Due problemi di incollamento}
\author{\small Michel Hua\thanks{Scritto con un assistente, Claude Opus~5.5 (Anthropic), poi riletto dall'autore.}}
\date{\small Novembre 2026}

\begin{document}
\maketitle
\vspace{-1.6em}

\noindent\emph{Introduzione.} Sono dati dei frammenti, ciascuno nella propria lingua. Li si
tratta come gli oggetti locali di un campo; ciò che abbozzano insieme è una sezione globale, e
ciò che risolvono si legge su una base. Le ostruzioni sono classi di coomologia, ed è queste che
si tratta di leggere.

\section{Il campo delle letture}

\begin{defi}
Sia $P$ un insieme finito di \emph{frammenti}, e $S$ un insieme di simboli. Una parte
$U \subset P$ è \emph{autonoma} se porta il contesto necessario per essere compresa da sola. Le
parti autonome e le loro inclusioni formano una categoria $C$, che si munisce della topologia
generata dalle famiglie $(U_i \subset U)$ di unione $U$. Un frammento isolato in generale non è
autonomo: è ciò che rende non banale questa topologia.
\end{defi}

\begin{defi}
Il \emph{campo delle letture} $\Lec$ associa a $U$ il gruppoide $\Lec(U)$ i cui oggetti sono le
letture di $U$ — un senso dato a ogni frammento, e un'applicazione che manda ogni simbolo usato
in $U$ sull'oggetto che designa — e le cui frecce $\ell \to \ell'$ sono le
\emph{ridenominazioni} dei simboli usati che trasformano $\ell$ in $\ell'$. La restrizione a
$V \subset U$ dimentica i frammenti fuori da $V$.
\end{defi}

\begin{prop}
$\Lec$ è un precampo: per $\ell, \ell' \in \Lec(U)$, il prefascio
$V \mapsto \mathrm{Isom}(\ell|_V, \ell'|_V)$ è un fascio. Un dato di discesa $(\ell_i, g_{ij})$
relativo a un ricoprimento $\U$ — letture $\ell_i \in \Lec(U_i)$ e ridenominazioni
$g_{ij} \colon \ell_j|_{U_{ij}} \to \ell_i|_{U_{ij}}$ tali che $g_{ij}g_{jk} = g_{ik}$ — è
effettivo, cioè proviene da una lettura di $U$ in una notazione unica, se e solo se la sua classe
è nulla in $\check H^1(\U, G)$, dove $G$ è il gruppo delle permutazioni di $S$.
\end{prop}

\begin{proof}
Ridenominazioni compatibili su un ricoprimento si incollano, poiché una ridenominazione non è che
un'applicazione definita simbolo per simbolo. Se $g_{ij} = h_i^{-1}h_j$, le letture $h_i\ell_i$
coincidono sulle intersezioni e si incollano; il viceversa è la definizione di banalità della
classe.
\end{proof}

Il campo associato $\tilde\Lec$ ha dunque più oggetti di $\Lec$: sono i \emph{testi}, dati di
discesa che si accetta di non ridurre a una notazione unica. Un oggetto di $\tilde\Lec(P)$ è un
testo di $P$; un oggetto di $\Lec(P)$ ne è una \emph{realizzazione}.

\begin{coro}
Se il nervo di $\U$ è un albero, ogni testo relativo a $\U$ ammette una realizzazione.
\end{coro}

\begin{proof}
$\check H^1$ di un albero è banale a coefficienti in un gruppo qualsiasi: si sceglie $h$ in una
radice e lo si propaga lungo i lati con $h_j = h_ig_{ij}$.
\end{proof}

\begin{prop}
Supponiamo la lettura \emph{localmente ben posta}: ogni frammento appartiene a una parte autonoma
che ammette una lettura, e due letture di una stessa parte sono localmente isomorfe. Allora
$\tilde\Lec$ è un gerbo; il suo legame è il fascio $\mathcal A$ degli automorfismi di una
lettura, le \emph{simmetrie di notazione}. Un testo di $P$ esiste se e solo se questo gerbo è
neutro; quando $\mathcal A$ è abeliano, l'ostruzione è una classe di $H^2(C, \mathcal A)$, e due
testi differiscono per un torsore sotto $\mathcal A$, classificato da $H^1(C, \mathcal A)$
(Giraud).
\end{prop}

\begin{proof}
Le due ipotesi sono, parola per parola, gli assiomi di un gerbo; il resto è la teoria di Giraud
dei gerbi a legame abeliano.
\end{proof}

\begin{rema}
Si vedono così due piani di ostruzione. Il primo ($H^2$) dice che non c'è alcun testo: le
simmetrie locali non si incollano. Il secondo ($H^1$) dice che un testo c'è, ma nessuna
realizzazione in una notazione unica. Una classe non nulla non è un errore: nasce da un ciclo di
frammenti in cui uno stesso nome designa due oggetti, e la si uccide ingrandendo $S$, non
correggendo una lettura.
\end{rema}

\begin{defi}
Un frammento può ammettere più letture. Sia $\Lambda(U) \subset \mathrm{Ob}\,\Lec(U)$ l'insieme
delle letture \emph{ammissibili} di $U$, e $E(\lambda)$ l'insieme degli enunciati che $\lambda$
afferma. Un enunciato è \emph{ricostruito} su $U$ se appartiene a
\[E(U) = \textstyle\bigcap_{\lambda \in \Lambda(U)} E(\lambda),\]
e \emph{corretto} se è vero così com'è scritto, eventualmente dopo l'aggiunta di un'ipotesi che
il frammento tace e che si dichiara.
\end{defi}

Restringere $\Lambda$ fa crescere $E$, e un enunciato ricostruito e corretto resta tale. Una
verifica \cache{formale} controlla dunque la matematica di una lettura, mai la sua fedeltà al
frammento: le due cose si misurano separatamente.

\section{Il punto di vista relativo}

Tutto ciò che precede si colloca al di sopra di una base: un insieme ordinato $T$, munito della
sua topologia di Alexandrov, che si chiamerà il \cache{tempo}. I frammenti vi sono collocati da un
morfismo $P \to T$; si indica con $t_P$ il limite superiore della sua immagine. Al di sopra di
$T$ vive un fascio crescente $t \mapsto C_{<t}$, gli enunciati che la comunità \cache{matematica}
riteneva acquisiti prima di $t$, e un insieme $Q$ di \cache{domande}, ciascuna con la sua data $t_q$. Un
ordine dei frammenti compatibile con gli indizi di posizione esiste se e solo se tali indizi sono
privi di cicli (Szpilrajn).

\begin{prop}
Sia $R \subset Q \times E$ la relazione « l'enunciato $e$ tocca la domanda $q$ ». Le applicazioni
$A \mapsto A'$ e $B \mapsto B'$ che essa definisce tra parti di $Q$ e di $E$ formano una
corrispondenza di Galois, e le coppie chiuse un reticolo completo.
\end{prop}

\begin{proof}
È la polarità di Birkhoff di una relazione binaria; le coppie chiuse sono i « concetti » di
Ganter e Wille.
\end{proof}

\begin{defi}
Sia $\vdash$ un sistema deduttivo in cui ogni domanda chiusa $q$ si enuncia con una formula
$\varphi_q$, e $\Cl X$ la chiusura deduttiva di $X$. Una \emph{risposta} a $q$ è una parte finita
$B \subset E(P)$ tale che $\varphi_q$ o $\neg\varphi_q$ appartenga a $\Cl B$; essa è
\emph{attestata} se ogni enunciato di $B$ porta il suo frammento, \emph{verificata} se la
deduzione è stata controllata, \emph{anticipata} se tutti i suoi frammenti precedono $t_q$.
\end{defi}

\begin{prop}
Se $\varphi_q \in \Cl{E(P)}$, esiste una risposta minimale rispetto all'inclusione, e l'insieme
delle domande risolte cresce con $P$.
\end{prop}

\begin{proof}
Una dimostrazione usa solo un numero finito di ipotesi.
\end{proof}

\begin{rema}
Una risposta anticipata è una sezione al di sopra dell'aperto $\{t < t_q\}$, che non contiene la
data della domanda: la domanda aveva una risposta prima di essere posta. Darla a posteriori non
stabilisce alcuna priorità; fornisce, con i frammenti, il materiale di una \cache{storia}. Una domanda
può anche avere risposta soltanto nel testo incollato, e in nessuna lettura locale.
\end{rema}

Sia infine $D(P) \subset E(P)$ l'insieme degli enunciati di cui i frammenti stessi danno una
dimostrazione.

\begin{conj}
Ogni enunciato corretto di $E(P)$ appartiene a
\[\Cl{D(P) \cup C_{<t_P}} .\]
Lo spazio degli enunciati dimostrabili a cui mirano i frammenti si ricostruisce a partire da essi
e da ciò che la comunità sapeva allora.
\end{conj}

\begin{rema}
Un controesempio sarebbe un enunciato corretto che i frammenti affermano senza dimostrarlo e che
nulla, al loro tempo, dimostrava: un'anticipazione. È cercando i controesempi che si trova ciò
che $P$ ha di nuovo; ogni risposta anticipata che cade fuori da $\Cl{D(P) \cup C_{<t_P}}$ ne è
uno.
\end{rema}

\begin{rema}[Riflessione]
Si sarebbe potuto attaccare a colpi di martello: fissare in anticipo una notazione e piegarvi ogni
frammento. Lo yoga qui è inverso. Si lascia ogni frammento nella sua lingua, si annota soltanto
come le lingue si traducono sui ricoprimenti, e si aspetta che la struttura parli: la coomologia
dice se un testo esiste e se ammette una realizzazione. Quando non si annulla, non dice che la
lettura è cattiva; dice dove guardare.
\end{rema}

\vspace{0.3em}
\noindent{\small\textbf{Bibliografia.}
G.~Birkhoff, \emph{Lattice Theory}, Amer. Math. Soc. Colloq. Publ.~25, 1940. ---
B.~Ganter, R.~Wille, \emph{Formal Concept Analysis. Mathematical Foundations}, Springer, 1999. ---
J.~Giraud, \emph{Cohomologie non abélienne}, Grundlehren~179, Springer, 1971. ---
R.~Godement, \emph{Topologie algébrique et théorie des faisceaux}, Hermann, 1958. ---
E.~Szpilrajn, « Sur l'extension de l'ordre partiel », \emph{Fund. Math.} 16 (1930), 386--389.}

\end{document}
